% example.tex - an example workterm report that can be used as a template for writing NE 
% workterm reports

% We begin by calling the workreport class which includes all the definitions for the macros
% and environments we will be using throughout the document.
\documentclass{workreport}

% We will use some packages to add some functionality to our script.
% The lipsum package is only used to create placeholder text for this example. It can be
% deleted when your own text has been filled into the sections.
\usepackage{lipsum}
% The amsmath package provides various features to facilitate writing math formulas.
\usepackage{amsmath}

%%%%%%%%%%%%%%%%%%
%  IMPORTANT  INFORMATION  %
%%%%%%%%%%%%%%%%%%

% Fill in the title of your report
\title{
	Example Work Term Report Style
}

% Fill in your name
\author{Regina Phalangey}

% Fill in your student number
\studentnumber{40440404}

% Fill in your student ID
\studentid{example}

% Fill in your previous academic term
\previousterm{2B}

% Fill in the name of your university
\university{University of Waterloo}

% Fill in your faculty
\faculty{Engineering}

% Fill in your program
\program{Nanotechnology Engineering}

% Fill in your program director
\programdirector{John Doe}

% Fill in the address of your university
\universityaddress{
	200 University Ave. West \\
	Waterloo, ON \\
	N2L 3G1
}

% Fill in your employers name
\employername{Employer Inc.}

% Fill in your employers address
\employeraddress{
	200 Payment St. \\
	Waterloo, ON \\
	A1A 1A1
}

% We will now begin writing the document
\begin{document}

%%%%%%%%%%%
% FRONT MATTER %
%%%%%%%%%%%

% The \begin{frontmatter} command will allow the use of the following environments:
% letter_of_submittal, contributions, summary, conclusions, recommendations, toctree
% (table of contents) and listoffiguresandtables (list of figures and list of tables).
\begin{frontmatter}

% The letter_of_submittal environment requires two inputs, the first being the greeting
% and the second being the closing.
\begin{letter_of_submittal}{Dear Dr. Doe:}{Sincerely,}

	% You're letter body.
	\lipsum[1-2]

\end{letter_of_submittal}

% The contribution environment only requires the contributions text.
\begin{contributions}

		\lipsum[1-3]

\end{contributions}

% The summary environment only requires the summary text.
\begin{summary}

		\lipsum[1-4]

\end{summary}

% The conclusions environment only requires the conclusions text.
\begin{conclusions}

		\lipsum[1-3]

\end{conclusions}

% The recommendations environment only requires the recmmendations text.
\begin{recommendations}

		\lipsum[1-3]

\end{recommendations}

% The following command generates a table of contents (toc) that automatically keeps
% track of sections and their page numbers.
\toctree

% The following command generates a list of figures and a list of tables that
% automatically keep track of properly inserted figures and tables. The commands for
% properly inserting figures and tables will be explained in the body section of this
% document.
\listoffiguresandtables

\end{frontmatter}

%%%%%%
% BODY  %
%%%%%%

% The \begin{body} command will create the body environment. This environment
% includes the formatting required for the body
\begin{body}

% The \section{title} command creates a section in the body of your report. You must
% specify the title of the section as the commands input. These sections headers will
% automatically be added to the table of contents.
\section{Introduction}

	\lipsum[1-6]

\section{Background}

	\lipsum[1-2]-

% The \subsection{title} command creates a subsection for the above mentioned
% section. You must specify the title of the section as the commands input. These
% subsection headers will automatically be added to the table of contents, below the
% proper section.
\subsection{A Subsection}

	\lipsum[1]

\subsection{A Second Subsection}

	\lipsum[2]

% The \subsubsection{title} command works exactly like the \section{title} and
% \subsection{title} commands above.
\subsubsection{A Subsection in a Subsection}

	\lipsum[1]

\subsubsection{Sub-ception}

	\lipsum[1]

	BWOOOONG!

\section{Data}

	\lipsum[1-2]

% This is where I stopped commenting (2016-04-13)
Here's a reference to table \ref{tbl:exampletable}

\begin{table}
	\centering
	\begin{tabular}{|c|c|c|} \hline
		title & number & decimal \\ \hline
		heading & 1 & 1.7 \\ \hline
		another heading & 2 & 3.4 \\ \hline
	\end{tabular}
	\caption{A table demonstrating \LaTeX \, formatting}
	\label{tbl:exampletable}
\end{table}

\section{Math}

\lipsum[1]

\subsection{The Bloch Equations}

	\begin{align}
		\frac{d\, M_x(t)}{dt} &= \gamma(\mathbf{M}(t) \times \mathbf{B}(t))_x - \frac{M_x(t)}{T_2} \\
		\frac{d\, M_y(t)}{dt} &= \gamma(\mathbf{M}(t) \times \mathbf{B}(t))_y - \frac{M_y(t)}{T_2} \\
		\frac{d\, M_z(t)}{dt} &= \gamma(\mathbf{M}(t) \times \mathbf{B}(t))_z - \frac{M_z(t) - M_0}{T_1}
	\end{align}

\subsection{The Schrodinger Equation}

	\begin{equation}
		i\hbar \frac{\partial \, \Psi}{\partial t} = \hat{H}\Psi
	\end{equation}

\section{Graphics}
	Here's a stock image of a computer. It comes with a CC0 license, enabling
	free distribution for commercial and personal use. It's a JPEG

	\begin{figure}[!ht]
		\centering
		\label{fig:stock_computer}
		\includegraphics[width=5cm]{./stock-image.jpg}
		\caption{A stock computer, I just wanted a picture here}
	\end{figure}

\section{Conclusions}
	\gls{LaTeX} is awesome.

\section{Recommendations}

	Use more \LaTeX. Also, use more \gls{Unix}. Also, citations like
	\cite{schweiger_principles_2001} will lead to a better quality of work.

\end{body}

\newglossaryentry{LaTeX}{
	name={\LaTeX},
	description={A document preparation system used for
		high-quality typesetting. It is most often used
		for medium-to-large technical or scientific documents.
		\LaTeX is not a word processor!
	}
}

\newglossaryentry{Unix}{
	name={Unix},
	description={A multi-user multi-tasking operating system
	developed by Ken Thompson and Dennis Ritchie in the 1970s, which has inspired many features behind Mac OSX and Linux
}
}

\addbibresource{./example.bib}


\begin{backmatter}

	\begin{appendix}{Here's An Appendix}
		\lipsum[1-2]

	\end{appendix}

	\begin{appendix}{Here's Another Appendix}
		\lipsum[1-2]
	\end{appendix}

	\begin{appendix}{A Figure in An Appendix}
		\lipsum[1-2]

		\begin{figure}[!ht]
			\centering
			\includegraphics[width=5cm]{./stock-image.jpg}
			\caption{It's back!}
		\end{figure}

	\end{appendix}

	\begin{appendix}{A Table in An Appendix}
		\lipsum[1-2]

		\begin{table}
			\centering
			\begin{tabular}{|c|c|c|} \hline
				\textbf{Data} & \textbf{Integer} & \textbf{Float} \\ \hline
				Foo & 1 & 1.1 \\ \hline
			\end{tabular}
			\caption{A Table in an Appendix, displaying the correct numbering}
		\end{table}

	\end{appendix}

\end{backmatter}

\end{document}
